\section{Conclusions}
\label{sec:conclusions}

A trained assistance dog's alert is only as useful as the human who perceives it. We built the missing half of an IMU-based alert collar---a streaming verification gate, an owner $\rightarrow$ caregiver $\rightarrow$ emergency escalation ladder, and a backend watchdog---and evaluated the whole chain with methods chosen to find its weaknesses. The escalation protocol can be made provably correct: model checking confirmed its safety properties, found and let us fix a watchdog design flaw, and led us to an integration bug no test had caught. The detector is the bottleneck. Scored on continuous recordings, it raises \R{cont.rfe.r60.fah} false alarms per hour of training-session motion; the escalation design can tolerate about \R{req.maxfah} per hour. Closing that gap---with all-day data, rotation-robust models, and dogs trained to repeat an unanswered alert---is the work that remains before a collar like this can be trusted by anyone.
